% Options for packages loaded elsewhere
\PassOptionsToPackage{unicode}{hyperref}
\PassOptionsToPackage{hyphens}{url}
%
\documentclass[
]{article}
\usepackage{amsmath,amssymb}
\usepackage{iftex}
\ifPDFTeX
  \usepackage[T1]{fontenc}
  \usepackage[utf8]{inputenc}
  \usepackage{textcomp} % provide euro and other symbols
\else % if luatex or xetex
  \usepackage{unicode-math} % this also loads fontspec
  \defaultfontfeatures{Scale=MatchLowercase}
  \defaultfontfeatures[\rmfamily]{Ligatures=TeX,Scale=1}
\fi
\usepackage{lmodern}
\ifPDFTeX\else
  % xetex/luatex font selection
\fi
% Use upquote if available, for straight quotes in verbatim environments
\IfFileExists{upquote.sty}{\usepackage{upquote}}{}
\IfFileExists{microtype.sty}{% use microtype if available
  \usepackage[]{microtype}
  \UseMicrotypeSet[protrusion]{basicmath} % disable protrusion for tt fonts
}{}
\makeatletter
\@ifundefined{KOMAClassName}{% if non-KOMA class
  \IfFileExists{parskip.sty}{%
    \usepackage{parskip}
  }{% else
    \setlength{\parindent}{0pt}
    \setlength{\parskip}{6pt plus 2pt minus 1pt}}
}{% if KOMA class
  \KOMAoptions{parskip=half}}
\makeatother
\usepackage{xcolor}
\usepackage{longtable,booktabs,array}
\usepackage{calc} % for calculating minipage widths
% Correct order of tables after \paragraph or \subparagraph
\usepackage{etoolbox}
\makeatletter
\patchcmd\longtable{\par}{\if@noskipsec\mbox{}\fi\par}{}{}
\makeatother
% Allow footnotes in longtable head/foot
\IfFileExists{footnotehyper.sty}{\usepackage{footnotehyper}}{\usepackage{footnote}}
\makesavenoteenv{longtable}
\usepackage{graphicx}
\makeatletter
\def\maxwidth{\ifdim\Gin@nat@width>\linewidth\linewidth\else\Gin@nat@width\fi}
\def\maxheight{\ifdim\Gin@nat@height>\textheight\textheight\else\Gin@nat@height\fi}
\makeatother
% Scale images if necessary, so that they will not overflow the page
% margins by default, and it is still possible to overwrite the defaults
% using explicit options in \includegraphics[width, height, ...]{}
\setkeys{Gin}{width=\maxwidth,height=\maxheight,keepaspectratio}
% Set default figure placement to htbp
\makeatletter
\def\fps@figure{htbp}
\makeatother
\setlength{\emergencystretch}{3em} % prevent overfull lines
\providecommand{\tightlist}{%
  \setlength{\itemsep}{0pt}\setlength{\parskip}{0pt}}
\setcounter{secnumdepth}{-\maxdimen} % remove section numbering
\ifLuaTeX
  \usepackage{selnolig}  % disable illegal ligatures
\fi
\IfFileExists{bookmark.sty}{\usepackage{bookmark}}{\usepackage{hyperref}}
\IfFileExists{xurl.sty}{\usepackage{xurl}}{} % add URL line breaks if available
\urlstyle{same}
\hypersetup{
  hidelinks,
  pdfcreator={LaTeX via pandoc}}

\author{}
\date{}

\begin{document}

\hypertarget{inter-channel-statistical-analysis-and-local-dispersion-indicators}{%
\section{Inter-Channel Statistical Analysis and Local Dispersion
Indicators}\label{inter-channel-statistical-analysis-and-local-dispersion-indicators}}

The empirical study of candidate pair distributions along the coordinate
axes of the \(\mathbb{R}_{30}\) modular structure reveals remarkable
geometric and regularity properties. By isolating the local behavior of
flows using rolling-window analyses and variance computations, it
becomes possible to quantify the robustness of the model against the
elimination mechanisms of the sieve.

\hypertarget{1-experimental-protocol-description}{%
\subsection{1. Experimental Protocol
Description}\label{1-experimental-protocol-description}}

The model segments the space of even solutions into distinct channels
derived from the multiplicative group
\((\mathbb{Z}/30\mathbb{Z})^\times\). For a given even integer \(N\),
the candidates \((p, q)\) satisfying the relation \(p + q = N\) are
grouped according to their residual congruence classes mod 30.

\hypertarget{figure-1-combinatorial-activation-matrix-modulo-30}{%
\subsubsection{Figure 1: Combinatorial Activation Matrix Modulo
30}\label{figure-1-combinatorial-activation-matrix-modulo-30}}

\includegraphics{/home/mycmon/Documents/7. TH1-TH11-C1-C5/TH12/fig1_activation_matrix.png}

This matrix maps the activation of finite residual classes and
illustrates how multiplicative interactions organize themselves for a
specific value of \(N\).

\begin{itemize}
\item
  XSG = Not Sophie Germain Number
\item
  SG = Sophie Germain Number
\end{itemize}

To finely measure dispersion and local regularity within these channels,
the space of each admissible channel is partitioned into \(B\)
consecutive sub-blocks (for instance, 8 successive blocks each
containing a fixed sample of 1,000 candidate pairs). A local sieve
process of level \(z = \sqrt{N}\) is applied to each block to evaluate
the survival rate of the pairs.

\hypertarget{figure-2-circular-projection-of-the-phase-circle--ux1d445-30}{%
\subsubsection{\texorpdfstring{Figure 2: Circular Projection of the
Phase Circle
\(\mathbb{R}_{30}\)}{Figure 2: Circular Projection of the Phase Circle \textbackslash mathbb\{R\}\_\{30\}}}\label{figure-2-circular-projection-of-the-phase-circle--ux1d445-30}}

\includegraphics{/home/mycmon/Documents/7. TH1-TH11-C1-C5/TH12/fig2_circle_R30.png}

Geometric visualization of the \(\mathbb{R}_{30}\) phase space,
highlighting the rigid symmetries of the 8 prime number classes.

The following statistical indicators are systematically computed for
each channel:

\begin{itemize}
\item
  \textbf{Local mean (\(\mu\)):} The average survival rate across all
  blocks.
\item
  \textbf{Standard deviation (\(\sigma\)) and Variance (\(\sigma^2\)):}
  Measuring the rate fluctuation from one block to another.
\item
  \textbf{Dispersion Index (or Fano Factor, \(I_D = \sigma^2 / \mu\)):}
  Characterizing the mathematical nature of the underlying statistical
  distribution.
\end{itemize}

\hypertarget{2-analysis-of-statistical-indicators-and-interpretation}{%
\subsection{2. Analysis of Statistical Indicators and
Interpretation}\label{2-analysis-of-statistical-indicators-and-interpretation}}

Data collected over several orders of magnitude (from \(10^8\) up to
\(5 \times 10^9\)) reveal three major characteristics of the modular
organization.

\hypertarget{figure-3-spectral-envelope-and-sieve-waves}{%
\subsubsection{Figure 3: Spectral Envelope and Sieve
Waves}\label{figure-3-spectral-envelope-and-sieve-waves}}

\includegraphics{/home/mycmon/Documents/7. TH1-TH11-C1-C5/TH12/fig3_spectral_envelope.png}

Illustration of the sieve\textquotesingle s exclusion forces and their
harmonic modulations at a macroscopic scale.

\hypertarget{a-extreme-under-dispersion-and-the-fano-factor}{%
\subsubsection{A. Extreme Under-Dispersion and the Fano
Factor}\label{a-extreme-under-dispersion-and-the-fano-factor}}

The most significant indicator lies in the systematic suppression of the
dispersion index (\(I_D\)). In a purely stochastic model---where the
distribution of prime numbers would follow a Poisson process
disconnected from any modular constraint---the dispersion index would
fluctuate around \(1.0\). However, empirical observations show uniformly
low values, frequently dropping below the \(0.1\) threshold.

\hypertarget{plot-b-dispersion-index-analysis-fano-factor-per-channel-vs-pure-randomness}{%
\subsubsection{Plot B: Dispersion Index Analysis (Fano Factor) per
Channel vs. Pure
Randomness}\label{plot-b-dispersion-index-analysis-fano-factor-per-channel-vs-pure-randomness}}

The chart below illustrates this structural gap: while pure randomness
sits on the red dashed line (\(1.0\)), the indexes of the actual
channels remain crushed near the bottom, confirming a strict
mathematical under-dispersion.

The alignment of sieve waves with the geometry of the 30-wheel generates
a smoothing effect: candidate eliminations do not cluster chaotically,
preventing the local emergence of large empty gaps or absolute "deserts"
within the channel.

\includegraphics{/home/mycmon/Documents/7. TH1-TH11-C1-C5/TH12/monfette_dispersion_1000000008.png}

\hypertarget{b-equidistribution-of-simultaneous-channels}{%
\subsubsection{B. Equidistribution of Simultaneous
Channels}\label{b-equidistribution-of-simultaneous-channels}}

Examining the survival means (\(\mu\)) among the different channels of
the same even target \(N\) shows remarkable homogeneity. At the scale of
\(N = 1\) Billion, the channels configured for this point deploy with
nearly identical survival averages: \textbf{7.12\%} for the \((1,17)\)
channel, \textbf{6.73\%} for the \((7,11)\) channel, and \textbf{7.02\%}
for the \((19,29)\) channel.

\hypertarget{plot-a-local-evolution-of-the-survival-rate-per-consecutive-block}{%
\subsubsection{Plot A: Local Evolution of the Survival Rate per
Consecutive
Block}\label{plot-a-local-evolution-of-the-survival-rate-per-consecutive-block}}

The plot below graphically displays the smooth and parallel trajectories
of these flows across successive candidate blocks, ruling out any local
anomalies or unexpected drops.

This observation validates the steady behavior of Dirichlet characters
associated with the residual classes of the 30-wheel. The exclusion
forces of the sieve are applied equitably, without introducing any major
distortion or bias in favor of a specific tunnel.

\includegraphics{/home/mycmon/Documents/7. TH1-TH11-C1-C5/TH12/monfette_taux_survie_1000000008.png}

\hypertarget{figure-4-the-goldbach-comet-and-its-macroscopic-distribution}{%
\subsubsection{Figure 4: The Goldbach Comet and its Macroscopic
Distribution}\label{figure-4-the-goldbach-comet-and-its-macroscopic-distribution}}

Classical macroscopic perspective of the pair distribution, serving as a
global baseline comparison for the microscopic regularity observed
within the channels.

\includegraphics{/home/mycmon/Documents/7. TH1-TH11-C1-C5/TH12/fig4_goldbach_comet.png}

\hypertarget{c-the-global-non-extinction-mechanism}{%
\subsubsection{C. The Global Non-Extinction
Mechanism}\label{c-the-global-non-extinction-mechanism}}

The asymptotic trajectory analyzed up to 5 Billion
(\(N = 5\,000\,000\,030\)) confirms the geometric robustness of the
framework. Despite a drastic increase in the number of tested prime
factors (moving from 1,226 to 7,001 factors), the standard deviation
(\(\sigma\)) remains confined to around \textbf{1\%}.

\hypertarget{figure-5-phase-duality-diagram-of-candidates}{%
\subsubsection{Figure 5: Phase Duality Diagram of
Candidates}\label{figure-5-phase-duality-diagram-of-candidates}}

Mapping of the symmetries and phase duality relations between the even
and odd components of the model.

Even within the sub-blocks most exposed to cross-eliminations, the
absolute number of surviving pairs remains significantly greater than
zero. The decay of variance as \(N\) grows renders the probability of a
simultaneous fluctuation across all admissible channels extremely
improbable.

\includegraphics{/home/mycmon/Documents/7. TH1-TH11-C1-C5/TH12/fig5_duality_diagram.png}

\hypertarget{figure-6-histogram-of-kappa--ux1d705--regularity-indexes}{%
\subsubsection{\texorpdfstring{Figure 6: Histogram of Kappa (\(\kappa\))
Regularity
Indexes}{Figure 6: Histogram of Kappa (\textbackslash kappa) Regularity Indexes}}\label{figure-6-histogram-of-kappa--ux1d705--regularity-indexes}}

Empirical statistical distribution validating the regularity indexes
measured across the entire computational campaign.

\includegraphics{/home/mycmon/Documents/7. TH1-TH11-C1-C5/TH12/fig6_kappa_histogram.png}

\hypertarget{3-summary-of-numerical-observations}{%
\subsection{3. Summary of Numerical
Observations}\label{3-summary-of-numerical-observations}}

The table below summarizes the evolution of the dispersion indicators
and confirms the stabilization of flows across the tested
configurations:

\begin{longtable}[]{@{}lllllll@{}}
\toprule\noalign{}
\textbf{Value of N} & \textbf{Sieve Level (z)} & \textbf{Number of
Primes} & \textbf{Modular Channel} & \textbf{Mean (μ)} &
\textbf{Standard Dev. (σ)} & \textbf{Dispersion Index (ID)} \\
\midrule\noalign{}
\endhead
\bottomrule\noalign{}
\endlastfoot
\textbf{\(100\,000\,018\)} & \(10\,000\) & 1,226 & \((11,17)\)
\((29,29)\) & \(6.812\%\) \(6.775\%\) & \(0.936\) \(0.772\) & \(0.1286\)
\textbf{\(0.0881\)} \\
\textbf{\(1\,000\,000\,008\)} & \(31\,622\) & 3,398 & \((1,17)\)
\((7,11)\) \((19,29)\) & \(7.125\%\) \(6.738\%\) \(7.025\%\) & \(1.036\)
\(0.583\) \(0.808\) & \(0.1508\) \textbf{\(0.0504\)} \(0.0931\) \\
\textbf{\(5\,000\,000\,030\)} & \(70\,710\) & 7,001 & \((1,19)\)
\((7,13)\) & \(5.225\%\) \(5.050\%\) & \(1.062\) \(0.886\) & \(0.2161\)
\(0.1554\) \\
\end{longtable}

\hypertarget{4-provisional-conclusion}{%
\subsection{4. Provisional Conclusion}\label{4-provisional-conclusion}}

The dispersion indicator approach shows that the simultaneous extinction
of channels is counteracted by two empirical factors: a very low
inter-block variance and a minimal Fano factor value. These observations
suggest that the modular geometric structure imposes an internal
regularity constraint, guaranteeing the maintenance of active candidate
flows for any sampled even \(N\).

\hypertarget{figure-7----3-uxd7-3-uxd7-3--topological-cube-of-minimal-configurations}{%
\subsubsection{\texorpdfstring{Figure 7: \(3 \times 3 \times 3\)
Topological Cube of Minimal
Configurations}{Figure 7: 3 \textbackslash times 3 \textbackslash times 3 Topological Cube of Minimal Configurations}}\label{figure-7----3-uxd7-3-uxd7-3--topological-cube-of-minimal-configurations}}

Three-dimensional modeling of non-empty overlap constraints (discussed
in the appendices or in connection with the Lean 4 formalization).

\includegraphics{/home/mycmon/Documents/7. TH1-TH11-C1-C5/TH12/fig7_cube_3x3x3.png}

\end{document}
